% ============================================================
%  CHƯƠNG 2: XÁC THỰC NGƯỜI DÙNG
%  Phụ trách: [Thành viên 2]
%  Packages: ui/splash, ui/auth, data/model/auth,
%            data/remote (AuthInterceptor, TokenAuthenticator),
%            utils/SessionManager
%  Ghi chú: File này được \input{} từ main.tex
% ============================================================

\chapter{Xác thực người dùng}

Chương này trình bày phân hệ xác thực của ứng dụng Trọ Tốt: đăng ký, đăng
nhập, điều hướng theo phiên và bảo vệ tài khoản bằng JWT. Nhóm viết toàn bộ
phân hệ bằng \tech{Java} theo kiến trúc MVVM ở Chương 1. Đây là lớp chắn đầu
tiên của ứng dụng: mọi API nghiệp vụ (đăng tin, chat, đánh giá) đều yêu cầu
token hợp lệ do phân hệ này cấp phát và duy trì.

% ============================================================
\section{Màn hình khởi động (Splash Screen)}
% ============================================================

\subsection{Yêu cầu chức năng}
Khi mở ứng dụng, người dùng không phải đăng nhập lại nếu phiên trước đó còn
hiệu lực. Màn hình khởi động hiển thị logo thương hiệu trong thời gian ngắn,
kiểm tra trạng thái phiên rồi tự chuyển đến màn hình phù hợp -- không yêu cầu
bất kỳ thao tác nào từ người dùng.

\subsection{Thiết kế giao diện}
Layout \code{activity\_splash.xml} tối giản: \code{FrameLayout} nền thương
hiệu, bên trong là \code{LinearLayout} căn giữa chứa \code{ImageView} (logo)
và \code{TextView} (tên ứng dụng). Màn hình chỉ hiển thị 1,5 giây nên không
đặt thêm thành phần tương tác.

\subsection{Chuyển hướng tự động (Auto-navigation)}

\subsection{Luồng kiểm tra phiên và chuyển hướng tự động}
\begin{enumerate}
    \item \code{SplashActivity} khởi chạy, Hilt inject sẵn
    \code{SessionManager} (annotation \code{@AndroidEntryPoint}).
    \item \code{Handler.postDelayed()} hẹn giờ 1,5 giây để logo hiển thị đủ
    lâu trước khi điều hướng.
    \item Hết thời gian chờ, Activity gọi \code{sessionManager.isLoggedIn()}
    -- giá trị đọc từ bộ nhớ phiên đã mã hóa (trình bày ở mục~\ref{sec:c2_session}).
    \item Nếu đã đăng nhập, ứng dụng mở thẳng \code{MainActivity}; ngược lại
    mở \code{AuthActivity} chứa luồng đăng nhập/đăng ký.
    \item \code{SplashActivity} gọi \code{finish()} để tự gỡ khỏi back stack
    -- nhấn Back từ màn hình kế tiếp sẽ không quay lại splash.
\end{enumerate}

\subsection{Triển khai trên Android}
\begin{itemize}
    \item \textbf{Lớp chính:} \code{SplashActivity} (package
    \code{ui/splash}), \code{SessionManager} (package \code{utils}).
    \item Không gọi API nào ở bước này: trạng thái phiên đọc cục bộ nên màn
    hình khởi động hoạt động cả khi mất mạng.
\end{itemize}

\textit{Đáp ứng CLO1: luồng khởi động hoàn chỉnh, tự động duy trì phiên đăng
nhập cho người dùng.}

% ============================================================
\section{Đăng ký tài khoản}
% ============================================================

\subsection{Yêu cầu chức năng}
Người dùng mới (người thuê hoặc chủ trọ) tạo tài khoản bằng số điện thoại và
email. Form đăng ký thu thập: họ tên, ngày sinh, giới tính, mật khẩu, nơi ở
hiện tại (tỉnh/thành phố $\rightarrow$ quận/huyện) và nghề nghiệp (sinh viên
/ đã đi làm) -- hai trường sau phục vụ mô hình gợi ý bài đăng. Ứng dụng chặn
đăng ký nếu người dùng chưa đủ 18 tuổi. Đăng ký thành công, ứng dụng đưa
người dùng về màn hình đăng nhập.

\subsection{Thiết kế giao diện}
Layout \code{fragment\_register.xml} đặt toàn bộ form trong \code{ScrollView}
để cuộn được trên màn hình nhỏ:
\begin{itemize}
    \item Cặp \code{TextInputLayout} + \code{TextInputEditText} (Material)
    cho số điện thoại, email, họ, tên, mật khẩu -- lỗi validation hiển thị
    ngay dưới từng ô nhập qua \code{setError()}.
    \item \code{AutoCompleteTextView} dạng dropdown cho giới tính, nghề
    nghiệp, tỉnh/thành phố và quận/huyện; danh sách quận/huyện nạp lại theo
    tỉnh đã chọn.
    \item Ô ngày sinh mở \code{DatePickerDialog}; ngày tối đa bị khóa ở mốc
    18 năm trước nên người dùng không thể chọn ngày sinh chưa đủ tuổi.
    \item \code{Button} đăng ký, \code{ProgressBar} trạng thái chờ,
    \code{MaterialCardView} (\code{cardError}) hiển thị lỗi từ server và
    \code{TextView} dẫn về màn hình đăng nhập.
\end{itemize}

% TODO: Bỏ comment sau khi chụp màn hình và lưu ảnh vào images/outputs/
% \begin{figure}[H]
%     \centering
%     \includegraphics[width=0.4\textwidth]{images/outputs/c2_register_screen.png}
%     \caption{Màn hình đăng ký tài khoản}
%     \label{fig:c2_register}
% \end{figure}

\subsection{Luồng xử lý và validation}
\begin{enumerate}
    \item \code{RegisterFragment} gom dữ liệu từ form, gọi
    \code{viewModel.register(...)}.
    \item \code{RegisterViewModel} xóa lỗi cũ rồi validate từng trường: số
    điện thoại theo mẫu \code{0} + 9 chữ số (regex), email theo regex, họ/tên
    tối thiểu 2 ký tự, mật khẩu tối thiểu 6 ký tự, ngày sinh phải đủ 18 tuổi
    (kiểm tra lại bằng \code{Calendar} dù UI đã khóa -- chặn cả trường hợp
    nhập tay).
    \item Trường nào sai, ViewModel phát thông báo qua \code{LiveData} lỗi
    tương ứng; Fragment observe và đẩy vào \code{TextInputLayout.setError()}.
    Luồng dừng tại đây, không gọi mạng.
    \item Hợp lệ, ViewModel phát \code{Resource.loading()} và tạo
    \code{RegisterRequest} gửi cho \code{AuthRepository}.
    \item Repository gọi API trên luồng IO qua RxJava, bọc kết quả vào
    \code{Resource} rồi trả về Main Thread.
    \item Thành công: Fragment hiện Toast \screen{Đăng ký thành công! Vui lòng
    đăng nhập} và điều hướng về \code{LoginFragment} qua Navigation Component.
    Thất bại: hiện \code{cardError} với thông báo từ server (ví dụ trùng
    email).
\end{enumerate}

\subsection{Triển khai trên Android}
\begin{itemize}
    \item \textbf{Lớp chính:} \code{RegisterFragment},
    \code{RegisterViewModel} (\code{@HiltViewModel}) thuộc package
    \code{ui/auth/register}; \code{AuthRepository} thuộc
    \code{data/repository}.
    \item \textbf{API sử dụng:} \code{POST api/auth/register} với body
    \code{RegisterRequest} $\rightarrow$ \code{RegisterResponse}.
    \item Danh sách tỉnh/quận đọc từ \code{LocationService} (dữ liệu hành
    chính đóng gói sẵn trong app), không tốn request mạng khi mở form.
\end{itemize}

\textit{Đáp ứng CLO1 và CLO2: form nhập liệu hoàn chỉnh với TextInputEditText,
AutoCompleteTextView, validation hai lớp và gọi API qua Retrofit.}

% ============================================================
\section{Đăng nhập}
% ============================================================

\subsection{Yêu cầu chức năng}
Người dùng đăng nhập bằng \textbf{email hoặc số điện thoại} (một ô nhập
chung -- identifier) cùng mật khẩu. Tùy chọn \screen{Ghi nhớ đăng nhập} lưu
lại identifier để điền sẵn ở lần mở sau. Đăng nhập thành công, ứng dụng vào
thẳng màn hình chính và duy trì phiên cho tới khi người dùng chủ động đăng
xuất hoặc refresh token hết hạn.

\subsection{Thiết kế giao diện}
Layout \code{fragment\_login.xml} gồm:
\begin{itemize}
    \item \code{ImageView} logo, \code{TextView} tên thương hiệu và tiêu đề.
    \item Hai cặp \code{TextInputLayout} + \code{TextInputEditText}: ô
    identifier (email/số điện thoại) và ô mật khẩu có nút ẩn/hiện.
    \item \code{MaterialCheckBox} (\code{cbRememberMe}) ghi nhớ đăng nhập;
    \code{TextView} \screen{Quên mật khẩu?} và liên kết chuyển sang đăng ký.
    \item \code{MaterialButton} (\code{btnLogin}) bị khóa trong lúc chờ,
    \code{ProgressBar} báo trạng thái và \code{MaterialCardView}
    (\code{cardError}) hiển thị lỗi đăng nhập sai.
\end{itemize}

% TODO: Bỏ comment sau khi chụp màn hình và lưu ảnh vào images/outputs/
% \begin{figure}[H]
%     \centering
%     \includegraphics[width=0.4\textwidth]{images/outputs/c2_login_screen.png}
%     \caption{Màn hình đăng nhập}
%     \label{fig:c2_login}
% \end{figure}

\subsection{Luồng xử lý}
\begin{enumerate}
    \item \code{LoginFragment} lấy identifier + mật khẩu, gọi
    \code{viewModel.login(...)}.
    \item \code{LoginViewModel} validate: identifier phải đúng định dạng
    email hoặc số điện thoại, mật khẩu tối thiểu 6 ký tự. Sai thì báo lỗi
    qua \code{LiveData}, không gọi mạng.
    \item ViewModel phát \code{Resource.loading()}; \code{AuthRepository}
    gửi \code{POST api/auth/login} trên luồng IO.
    \item Server trả về cặp token (access + refresh) và thông tin tài khoản;
    Repository lưu ngay vào \code{SessionManager} rồi phát
    \code{Resource.success()}.
    \item Nếu bật ghi nhớ, ViewModel lưu identifier để điền sẵn lần sau;
    đồng thời đăng ký FCM token lên server để nhận thông báo đẩy trên
    thiết bị này.
    \item Fragment điều hướng sang \code{MainActivity} với cờ
    \code{FLAG\_ACTIVITY\_CLEAR\_TASK} -- xóa sạch back stack, nhấn Back
    không quay lại được màn hình đăng nhập.
    \item Đăng nhập sai: \code{cardError} hiện thông báo lỗi do
    \code{ErrorHandler} chuẩn hóa từ phản hồi server.
\end{enumerate}

\subsection{Triển khai trên Android}
\begin{itemize}
    \item \textbf{Lớp chính:} \code{AuthActivity} (host Navigation),
    \code{LoginFragment}, \code{LoginViewModel} thuộc package
    \code{ui/auth/login}; \code{AuthRepository} thuộc \code{data/repository}.
    \item \textbf{API sử dụng:} \code{POST api/auth/login} với body
    \code{LoginRequest\{identifier, password\}} $\rightarrow$
    \code{LoginResponse} chứa \code{Token} và \code{Account}.
    \item Ba màn hình đăng nhập/đăng ký/quên mật khẩu là Fragment trong cùng
    \code{AuthActivity}, chuyển cảnh bằng \tech{Navigation Component} --
    không tạo Activity riêng cho từng bước.
\end{itemize}

\textit{Đáp ứng CLO2: chức năng login làm cơ sở phân quyền người thuê / chủ
trọ trong toàn ứng dụng.}

% ============================================================
\section{Quản lý phiên và bảo mật token}
\label{sec:c2_session}
% ============================================================

\subsection{Yêu cầu chức năng}
Phiên đăng nhập phải thỏa ba ràng buộc: token không bị lộ khi thiết bị bị
trích xuất dữ liệu, mọi request nghiệp vụ tự động mang token mà tầng UI không
phải bận tâm, và access token hết hạn phải được làm mới ngầm -- người dùng
không bị văng ra giữa chừng.

\subsection{Cơ chế JWT và lưu trữ token an toàn}
Server cấp hai JWT khi đăng nhập: \textbf{access token} (thời hạn ngắn, gắn
vào từng request) và \textbf{refresh token} (thời hạn dài, chỉ dùng để xin
access token mới). \code{SessionManager} lưu cả hai bằng
\code{EncryptedSharedPreferences} của thư viện AndroidX Security: khóa mã hóa
theo scheme \code{AES256\_SIV}, giá trị theo \code{AES256\_GCM}, master key
sinh và giữ trong Android Keystore. Dữ liệu phiên vì vậy vẫn an toàn kể cả
khi file preferences bị sao chép ra ngoài. Ngoài ra, \code{SessionManager}
còn bóc tách payload của JWT (Base64) để đọc \code{accountId} dùng cho kênh
chat/video call, thay vì tin vào dữ liệu rời rạc có thể thiếu.

\subsection{Tự động gắn token vào request}
\code{AuthInterceptor} (OkHttp Interceptor) chặn mọi request đi ra:
\begin{itemize}
    \item Gắn header \code{Authorization: Bearer <accessToken>} nếu phiên
    đang có token.
    \item Bỏ qua các đường dẫn \code{/auth/login}, \code{/auth/register},
    \code{/auth/refresh-token} -- tránh gửi kèm token cũ đã hết hạn vào
    chính các API xác thực.
    \item Bổ sung các header chung (\code{Accept},
    \code{Content-Type: application/json}) cho request có body.
\end{itemize}
Nhờ đặt ở tầng OkHttp, toàn bộ Repository gọi API mà không lặp lại mã xử lý
token ở bất kỳ đâu.

\subsection{Tự động làm mới token (Token Refresh)}
\code{TokenAuthenticator} cài đặt \code{okhttp3.Authenticator} -- OkHttp tự
kích hoạt khi server trả \code{401 Unauthorized}:
\begin{enumerate}
    \item Đếm số lần request đã thử lại; quá 2 lần thì bỏ cuộc, tránh lặp
    vô hạn.
    \item Nếu chính request refresh bị 401 (refresh token hết hạn) hoặc body
    lỗi báo \code{INVALID\_ACCESS\_TOKEN} (token bị can thiệp), xóa phiên
    ngay -- người dùng phải đăng nhập lại.
    \item Trường hợp còn lại, gọi đồng bộ \code{POST api/auth/refresh-token}
    với refresh token hiện có.
    \item Server chỉ trả access token mới; \code{SessionManager} lưu token
    mới nhưng \textbf{giữ nguyên refresh token cũ}.
    \item Trả về request gốc đã thay header Authorization -- OkHttp tự gửi
    lại, tầng UI không nhận biết có gián đoạn.
\end{enumerate}

\subsection{Triển khai trên Android}
\begin{itemize}
    \item \textbf{Lớp chính:} \code{SessionManager} (package \code{utils});
    \code{AuthInterceptor}, \code{TokenAuthenticator} (package
    \code{data/remote}) -- cả hai đăng ký vào \code{OkHttpClient} qua
    \code{NetworkModule} (Hilt).
    \item \textbf{API sử dụng:} \code{POST api/auth/refresh-token} với body
    \code{RefreshTokenRequest\{refreshToken\}} $\rightarrow$
    \code{RefreshTokenResponse\{accessToken\}}. Riêng API này khai báo bằng
    \code{retrofit2.Call} thay vì RxJava vì \code{Authenticator} chạy đồng
    bộ.
    \item Đăng xuất: xóa phiên trong \code{SessionManager} và hủy đăng ký
    FCM token trên server để thiết bị ngừng nhận thông báo.
\end{itemize}

\textit{Đáp ứng CLO2: bảo vệ tài khoản cá nhân và quyền riêng tư -- token mã
hóa khi lưu trữ, tự động làm mới phiên, thu hồi khi đăng xuất.}
